\documentclass[11pt]{article}
\usepackage[T1]{fontenc}
\usepackage{lmodern}
\usepackage[margin=1in]{geometry}
\usepackage{amsmath,amssymb,mathtools,microtype}
\usepackage{xurl}
\usepackage[colorlinks=true,linkcolor=blue,urlcolor=blue]{hyperref}
\hypersetup{pdftitle={The irrationality exponent of positive quadratic logarithms is at most 4},pdfauthor={Ryan Matthew Casper}}
\setlength{\emergencystretch}{3em}
\setcounter{secnumdepth}{0}
\providecommand{\tightlist}{\setlength{\itemsep}{0pt}\setlength{\parskip}{0pt}}
\title{The irrationality exponent of\\positive quadratic logarithms is at most 4}
\author{Ryan Matthew Casper}
\date{October 8, 2026}
\begin{document}
\maketitle
\section{Abstract}\label{abstract}

For every positive real algebraic number of degree two other than one,
we prove that its logarithm has irrationality exponent at most four. The
argument extends weighted logarithmic interpolation and determinant
estimates to the original quadratic field. Integer clearing, the field
norm, and a separate bound at the other embedding accommodate
nonintegral inputs and negative conjugates. The denominator threshold is
not made effective. A Lean formalization accompanies the proof.

\section{1. Target, conventions, and algebraic
setup}\label{target-conventions-and-algebraic-setup}

Let \(a\in\mathbb R\) satisfy \(a>0\), \(a\ne1\), and
\([\mathbb Q(a):\mathbb Q]=2\). Put \(\xi=\log a\). The target is \[
\forall\nu>4\;\exists Q\in\mathbb N\;\forall p\in\mathbb Z\;\forall q\in\mathbb N,
\quad q\ge Q\Longrightarrow
\left|\xi-\frac pq\right|\ge q^{-\nu}.                 \tag{1}
\] We may and will require \(Q\ge2\). The threshold may depend on
\(a,\nu\); no effectivity or uniformity in \(a\) is asserted. Fractions
need not be reduced. With the usual irrationality exponent defined using
errors \(|\xi-p/q|<q^{-\nu}\), (1) implies \(\mu(\xi)\le4\).

Write \(L=\mathbb Q(a)\subset\mathbb R\). Its other embedding is real:
if the monic minimal polynomial is \(X^2+tX+u\), its other root is
\(a'=-t-a\in\mathbb R\). Irreducibility makes \(a'\ne a\) and gives the
nontrivial field automorphism \(\tau:L\to L\), \(\tau(a)=a'\). No
positivity assumption is made on \(a'\). Choose an integer \(d\ge1\) for
which \(da\) is an algebraic integer. Such a choice follows, for
example, by taking the absolute value of a nonzero leading coefficient
of an integer polynomial for \(a\).

For a nonzero algebraic integer \(z\in L\), \[
z\tau(z)=N_{L/\mathbb Q}(z)\in\mathbb Z\setminus\{0\},
\qquad |z|\,|\tau(z)|\ge1.                           \tag{2}
\] These are standard quadratic-field facts. They are not consequences
of an absolute-value lower bound at either embedding separately.

We argue by contradiction for a fixed \(\nu>4\). If (1) fails, there are
approximants with arbitrarily large denominators and \[
|\xi-p/q|\le q^{-\nu}.                              \tag{3}
\] Set \(R=|\xi|+1\) and \(\rho=100\max(1,|\xi|)\). Every chosen
approximant with \(q\ge2\) has \(|p/q|\le R\). Since \(\xi\ne0\),
sufficiently large denominators also ensure \(p\ne0\), if that
convenient extra condition is required by a wrapper.

\section{2. Exact inherited inputs}\label{exact-inherited-inputs}

All source paths in this section are relative to the pinned
rational-logarithm dependency project. The namespace prefix \path{OAI}
is suppressed below where unambiguous. These statements have no
hypothesis that the multiplicative centres or the target logarithm are
rational.

\subsection{2.1 Interpolation}\label{interpolation}

Fix integers \(m,K\ge1\), positive rational degree weights
\(W=(w_0,w_1,\ldots,w_m)\), positive rational jet weights
\(V=(v_0,w_1/\theta,\ldots,w_m/\theta)\), \(0<\theta<1\), and rational
\(\sigma>0\). Require \[
K(1+3\sigma)^{m+1}\frac{\prod W_i}{\prod V_i}<1,
\qquad (1+\sigma)\theta<1,                           \tag{4}
\] and the separated-product inequalities in
\path{Logarithm.GeometryData.curve_separated_weight_products}: for
equal-size subsets \(I,J\subseteq\{0,\ldots,m\}\), if their largest
differing index is positive and belongs to \(I\), \[
\mathcal C(m,\sigma)\prod_{j\in J}V_j<\prod_{i\in I}W_i,
\quad
\mathcal C(m,\sigma)=
\operatorname{comparisonConstant}(m,\sigma).
                                                               \tag{5}
\] We use the actual source constant; only its finiteness and
independence from the centres matter here.

For arbitrary distinct nonzero \(y_j\in\mathbb C\) and arbitrary
\(c_{ji}\in\mathbb C\), the coefficient packets \[
[t^su^\beta]\,P\bigl(y_j(1+t),c_{j1}+u_1+\log(1+t),\ldots,
c_{jm}+u_m+\log(1+t)\bigr),\quad v_0s+w\beta/\theta<H,
\] can all be prescribed by polynomials of \(W\)-degree at most \(H\),
for all sufficiently large \(H=n\mathcal R\). The positive rational
scale \(\mathcal R\) depends only on the weights; the eventual threshold
may depend on the fixed centres.

\textbf{Declaration:}
\path{Logarithm.eventually_weighted_logarithmic_interpolation} in
\path{Logarithm/Interpolation.lean}, applied to the concrete
\path{GeometryData}. Its conclusion is the surjectivity of
\path{GeometricInterpolation.packetMap};
\path{GeometricInterpolation.WeightedPolynomials} supplies the degree
bound. Replacing each logarithm by a truncation \(G_i\) is allowed when
\(T_iv_0\ge w_i/\theta\), by
\path{eventually_weighted_truncated_interpolation}. The conversion to
the literal matrix is
\path{Logarithm.ScaledMatrix.truncatedLogMatrix_surjective_of_formalLog_packets}
in \path{ScaledMatrix.lean}.

The substantial geometric proof---curve inequality, blowup ampleness,
eventual jet restriction, and polynomial section degree bound---is
inherited through these declarations. This manuscript does not replace
it by a dimension count.

\subsection{2.2 Generic analytic determinant
ingredients}\label{generic-analytic-determinant-ingredients}

The following declarations are used with \(\omega=\xi\in\mathbb C\),
rational additive centres \(r_i=p_i/q_i\), and the entry factor
\(a^{jh}\).

\begin{itemize}
\tightlist
\item
  Finite row expansion with the centre factor retained:
  \path{Logarithm.ScaledTranslation.det_matrix_translation},
  \path{Logarithm/ScaledTranslation.lean}
\item
  Error \(j(r_i-\xi)\) in exponential weight form:
  \path{Logarithm.ScaledApproximationErrors.actual_error_exp},
  \path{Logarithm/ScaledApproximationErrors.lean}
\item
  Translation scalar estimate:
  \path{PiExponent.MatrixTranslationBounds.norm_rowScalar_exp_weight_difference},
  \path{OAI/NumberTheory/PiExponent/Approximation/MatrixTranslationBounds.lean}
\item
  Generic collision estimate:
  \path{Logarithm.ScaledCollision.formal_scaled_collision_bound},
  \path{Logarithm/ScaledCollision.lean}
\item
  Conversion of formal period coefficients to analytic row tests:
  \path{Logarithm.ScaledPeriodAnalytic.scaled_periodMonomial_coeff_eq_rowTest},
  \path{Logarithm/ScaledPeriodAnalytic.lean}
\item
  Two alternatives and finite summation:
  \path{PiExponent.DeterminantAnalyticBound.translated_summand_bound}
  and \path{log_norm_sum_le},
  \path{OAI/NumberTheory/PiExponent/Analysis/DeterminantAnalyticBound.lean}
\item
  Row and low-index lattice asymptotics:
  \path{PiExponent.MatrixCounting.tendsto_rowCount_normalized},
  \path{tendsto_lowIndexCount_normalized},
  \path{tendsto_collisionRatio},
  \path{OAI/NumberTheory/PiExponent/Approximation/MatrixCounting.lean}
\end{itemize}

The rational wrappers \path{LiteralAnalytic.actual_minor_analytic_bound}
and \path{GlobalMatrixInterpolation.cofinal_actualMatrix} \textbf{cannot
be applied directly} to \(a\). Sections 3 and 5 below reassemble their
generic ingredients and discharge the hypotheses. No new analytic
theorem is being silently assumed.

\section{3. Parameters are chosen before
approximants}\label{parameters-are-chosen-before-approximants}

Let \(\Lambda=4+\log4\), \(c=\log2/4\). Choose rational \(\theta,A,B,C\)
and real \(\eta\) satisfying \[
0<\theta<A<B<1,\quad C>1,\quad CB<1,\quad B<C\theta<1,
\quad0<\eta<1,
\] \[
g:=\nu\bigl(A(1-\eta)-\theta\bigr)-2(1-\theta)>0.     \tag{6}
\] This follows directly from
\texttt{PiExponent.exists\_parameters (nu/2)} since \(\nu/2>2\):
multiply its final gap by two. For an explicit elementary construction,
choose rational \(b\in(1/2,1-2/\nu)\), put \(\theta=1-s\), \(A=1-bs\)
with sufficiently small rational \(s>0\), use \(A^2<\theta\) to choose
\(C\in(A/\theta,1/A)\), and then choose \(A<B<\min(C\theta,1/C)\). A
sufficiently small positive \(\eta\) preserves the strict gap.

Fix \(0<\varepsilon<\min(g,1)\). Choose \(F>2/\theta\) such that
\(\nu/F<\varepsilon/3\). Set \[
K=\lfloor C^m\rfloor,\quad w_0=B^{-m},\quad
v_0=2K\theta^mw_0,
\quad L'=\log\max(1,|a'|),\quad \delta=\log d.        \tag{7}
\] All \(w_0,v_0\) are positive rational numbers, and \[
K(w_0/v_0)\theta^m=1/2,\qquad K\theta^m<1.
\] As \(m\to\infty\), \(K/w_0\to0\), \(1/w_0\to0\), \(m/v_0\to0\), while
\[
c_\infty:=\frac{c\eta^2K\theta^m}{(m+1)v_0A^m}
=\frac{c\eta^2(B/A)^m}{2(m+1)}\longrightarrow\infty. \tag{8}
\] The floor in \(K\) causes no problem: \(\lfloor C^m\rfloor/C^m\to1\).
These are also the identities underlying \path{exists_dimension_margin}
in the inherited parameter machinery.

Choose one finite \(m\ge1\) for which \(c_\infty>3\) and \[
\frac{2\Lambda Fm+3\log2}{v_0}
+\frac{(\rho+2\delta+L')K+\log(3/2)}{w_0}
<\varepsilon/3.                                    \tag{9}
\] Freeze \(m,K,w_0,v_0\). Choose rational \(\sigma>0\) sufficiently
small that (4) holds; one may also impose the inherited optional
condition \((1+3\sigma)^mK\theta^m<1\).

Only now choose \(m\) approximants from (3), successively, with
\(q_i\ge2\) and integer weights \(w_i=\lceil\log q_i\rceil\).
Arbitrarily large denominators allow the \(w_i\) to be as large as
desired. They can satisfy (5): cancel common factors in each product
comparison and solve for its largest differing weight, whose index
belongs to the degree-weight side; all remaining factors have smaller
indices and are already fixed. There are finitely many such comparisons
at each step. These thresholds depend on previously chosen weights, not
on the centre locations.

Also make \(w_*:=\min_iw_i\) large enough that \[
2\Lambda\sum_i\frac1{w_i}
+\frac{2\theta+\log4+\log(2K)+\nu+\log(2\rho K)+\log(KR+2)}{w_*}
<\varepsilon/3.                                    \tag{10}
\] This is possible because the finite numerator in (10) was fixed
before choosing the approximants, and \(\sum_i1/w_i\le m/w_*\).

Finally set \[
T_i=\lceil Fw_i/v_0\rceil,\qquad
G_i(t)=\sum_{1\le k<T_i}\frac{(-1)^{k+1}t^k}{k}.
\] Thus \(T_iv_0\ge Fw_i>w_i/\theta\). All parameters, approximants,
truncations, and centres are now frozen. The only remaining limit is
\(H\to\infty\).

\section{4. One selected algebraic minor and its arithmetic
bound}\label{one-selected-algebraic-minor-and-its-arithmetic-bound}

Take \(y_j=a^j\), \(c_{ji}=jr_i\), \(0\le j<K\). The \(y_j\) are
distinct and nonzero because \(a>0\), \(a\ne1\). The generic
interpolation input therefore gives full row rank at every sufficiently
large admissible height \(H=n\mathcal R\).

Rows are \((j,s,\beta)\) with \(v_0s+w\beta/\theta<H\); columns are
\((h,\gamma)\) with \(w_0h+w\gamma\le H\). All indices are nonnegative
integers. The entry is \[
E_{(j,s,\beta),(h,\gamma)}
=[t^su^\beta]\;a^{jh}(1+t)^h\prod_i(jr_i+G_i(t)+u_i)^{\gamma_i}\in L. \tag{11}
\] Choose a full-row square minor \(\Delta_H\ne0\), with \(M_H\) rows.
This selection is fixed in every estimate at that height. Its conjugate
is the same selection with \(a\) replaced by \(a'\); it is nonzero
because \(\tau\) is injective.

Define \[
b_H=\frac{\sum_{\rm rows}w\beta}{M_HH},\qquad0\le b_H\le\theta.
\] There is at least one row for \(H>0\), so \(M_HH>0\). Let
\(L_i=\operatorname{lcm}(1,\ldots,T_i)\) and \[
D_H=\prod_iL_i^{\lfloor H/w_i\rfloor}.
\] Multiply each entry by \(D_H\), its row factor
\(\prod_iq_i^{-\beta_i}\), and its column factor
\(d^{Kh}\prod_iq_i^{\gamma_i}\). If \(\beta\not\le\gamma\), the entry is
zero. Otherwise it becomes \[
D_H(da)^{jh}d^{(K-j)h}\binom\gamma\beta
[t^s](1+t)^h\prod_i(jp_i+q_iG_i(t))^{\gamma_i-\beta_i}. \tag{12}
\] Every coefficient of \(L_iG_i\) is an integer. Since
\(\gamma_i-\beta_i\le\gamma_i\le\lfloor H/w_i\rfloor\), the factor
\(D_H\) clears all polynomial coefficient denominators in (12). Thus
(12) is an algebraic integer in \(L\), not merely in \(L(i)\). The
determinant of the scaled matrix is \[
Z_H=C_H\Delta_H,\quad
C_H=D_H^{M_H}
\prod_{\rm selected\ columns}\left(d^{Kh}\prod_iq_i^{\gamma_i}\right)
\prod_{\rm rows}\prod_iq_i^{-\beta_i}>0.              \tag{13}
\] The positive rational number \(C_H\) need not be integral. The scaled
entries and their determinant are integral, which is what (2) needs.
Consequently \[
1\le C_H^2|\Delta_H|\,|\tau\Delta_H|.               \tag{14}
\]

For completeness, the clearing cost is \[
\frac{\log C_H}{M_HH}\le1-b_H+E_{\rm ar},
\quad E_{\rm ar}=
\frac{\Lambda Fm}{v_0}+\Lambda\sum_i\frac1{w_i}
+\frac\theta{w_*}+\frac{K\delta}{w_0}.               \tag{15}
\] Indeed \(\log L_i\le\Lambda T_i\), \(T_i\le Fw_i/v_0+1\), and
\(\lfloor H/w_i\rfloor\le H/w_i\), giving the first two terms. A
selected column satisfies
\(\sum_i\gamma_i\log q_i\le\sum_i\gamma_iw_i\le H\), and \(h\le H/w_0\),
giving \(1+K\delta/w_0\). For each row, \(\log q_i\ge w_i-1\) and
\(\sum_i\beta_i\le w\beta/w_*\le\theta H/w_*\). Thus its subtracted cost
is at least \(w\beta-\theta H/w_*\), proving (15). This exact row
subtraction is preserved twice in (14).

The lcm estimate is \path{PiExponent.Arithmetic.log_lcmUpto_le}; the
column and row inequalities are \path{column_log_cost_le} and
\path{row_log_cost_lower} in \path{Approximation/Arithmetic.lean}, or
their matrix wrappers. The new ordinary-integer version of (12) is the
isolated module \path{AlgebraicLog/IntegerClearing.lean}; the inherited
Gaussian clearing endpoint is deliberately not used to infer integrality
in \(L\).

\section{5. Distinguished analytic bound, with all
substitutions}\label{distinguished-analytic-bound-with-all-substitutions}

Put \[
f_{b,h,\gamma}(z)=[u^b]e^{hz}\prod_i(z+u_i)^{\gamma_i}
=\binom\gamma b e^{hz}z^{|\gamma|-|b|}
\] when \(b\le\gamma\), and zero otherwise. These are entire functions.
At row centre \(j\), write \(z_j(t)=j\xi+\log(1+t)\). Since
\(e^{\xi}=a\), \(e^{z_j(t)}=a^j(1+t)\) for \(|t|<1\). The additive
coordinates in (11) are \[
z_j(t)+u_i+e_{ji}+\tau_i(t),\qquad
e_{ji}=j(r_i-\xi),\quad \tau_i=G_i-\log(1+t).
\] The finite row expansion in the source list above therefore applies
literally. A row choice consists of a transverse multi-index \(b\) with
\(wb\le H\), a distribution of error/tail factors, and a coefficient
order \(k\le s\). The scalar is independent of the selected column.
Nonzero scalar terms have \(b\ge\beta\).

From (3) and \(w_i-1\le\log q_i\), \(|e_{ji}|\le Ke^\nu e^{-\nu w_i}\).
The tails begin in order \(T_i\). Apply
\path{norm_rowScalar_exp_weight_difference} with its error constant
\(\log(2K)+\nu\), weights \(w_i\), scale \(F\), lower weight \(w_*\),
and \(k\le s<H/v_0\). Its hypotheses hold because \(T_i\ge1\),
\(Fw_i\le v_0T_i\), and the approximation bound just proved. It gives
the scalar bound \[
\exp\{-\nu(wb-w\beta)+H E_{\rm tr}\},\qquad
E_{\rm tr}=\frac\nu F+\frac{\log2}{v_0}
+\frac{\log4+\log(2K)+\nu}{w_*}.                   \tag{16}
\]

Apply \path{formal_scaled_collision_bound} with row label \(j\), row
order \(\ell=s-k\), transverse grouping \(b\), column exponents
\(h,\gamma\), \(\omega=\xi\), and \(\rho\) chosen above. All its
hypotheses follow from \(0\le j<K\), the column degree budget,
\(\sum\ell\le M_HH/v_0\), \(\rho\ge100\), and \(\rho\ge100|\xi|\). Its
centre factor \(e^{j\xi h}\) equals exactly \(a^{jh}\). For \(n_b\) rows
in each group, the resulting determinant is at most \[
\exp\left\{-c\sum_bn_b^2+M_HH(E_{\rm hol}+o(1))\right\},
\quad E_{\rm hol}=\frac{\rho K}{w_0}+\frac{\log2}{v_0}
+\frac{\log(2\rho K)}{w_*}.                         \tag{17}
\] The source writes the last logarithm as
\(\log(200K)+\log(\rho/100)\), which is equal. The \(o(1)\) is
\texttt{Collision.collisionRemainder M\_H H}; its convergence follows
from the positive row-count asymptotic below and
\path{tendsto_collisionRemainder_of_normalized_pow}.

Let \(N_A(H)=\#\{b:wb\le AH\}\) and \[
c_H=\frac{c\eta^2M_H}{H N_A(H)}.
\] If at least \(\eta M_H\) rows use low indices, Cauchy--Schwarz gives
\(\sum n_b^2\ge\eta^2M_H^2/N_A(H)\). Otherwise the total index weight is
at least \(A(1-\eta)M_HH\). Also \(\sum wb\ge\sum w\beta=M_HHb_H\) for
every nonzero scalar product. Applying \path{translated_summand_bound}
gives exactly these two alternatives.

There are at most
\(Q_H=(\lfloor H\rfloor+1)^{2m}(\lfloor H/v_0\rfloor+1)\) choices per
row, as in \path{MatrixTranslation.row_term_count}; \(w_i\ge1\) and
\(s\le H/v_0\) discharge its hypotheses. Sum the at most \(Q_H^{M_H}\)
terms using \path{log_norm_sum_le}. Since \(\log Q_H/H\to0\), we obtain,
uniformly in the selected columns, \[
\frac{\log|\Delta_H|}{M_HH}\le
E_{\rm an}+o(1)+\max\{-c_H,-\nu(A(1-\eta)-b_H)\},
\quad E_{\rm an}=E_{\rm tr}+E_{\rm hol}.             \tag{18}
\]

The rational positive weights and \(q_i\ge2\) allow the lattice-counting
declarations in the table to be applied, giving \[
M_H\sim\frac{K\theta^mH^{m+1}}{(m+1)!v_0\prod_iw_i},
\quad N_A(H)\sim\frac{A^mH^m}{m!\prod_iw_i},
\quad c_H\longrightarrow c_\infty.                 \tag{19}
\] Thus every normalized error in (18) tends to zero after all other
data have been fixed.

\section{6. The other embedding: a new coarse
bound}\label{the-other-embedding-a-new-coarse-bound}

Apply the same coefficient formula (11) with \(a'\) in place of \(a\).
On \(|t|=1/2\), \(|G_i(t)|\le\sum_{k\ge1}2^{-k}/k=\log2\). On the
additional circles \(|u_i|=1\), \[
|jr_i+G_i(t)+u_i|\le KR+\log2+1<KR+2.
\] Repeated one-variable Cauchy coefficient bounds therefore give \[
|\tau E_{(j,s,\beta),(h,\gamma)}|
\le 2^s\max(1,|a'|)^{Kh}(3/2)^h(KR+2)^{|\gamma|}.
\] The \(u\)-circle radii equal one, so no additional
\(\beta\)-dependent cost appears. This estimate already includes all
binomial coefficients. Use \(s<H/v_0\), \(h\le H/w_0\), and
\(|\gamma|\le H/w_*\) to bound every entry by \(e^{HU}\), where \[
U=\frac{\log2}{v_0}+\frac{KL'+\log(3/2)}{w_0}
+\frac{\log(KR+2)}{w_*}.                            \tag{20}
\] The permutation formula for the determinant then gives \[
\frac{\log|\tau\Delta_H|}{M_HH}\le
U+\frac{\log(M_H!)}{M_HH}=U+o(1).                   \tag{21}
\] Indeed \(M_H=O(H^{m+1})\), so \(\log(M_H!)/(M_HH)\le\log M_H/H\to0\).
This proof is uniform in the selected minor. It uses neither
\(\log a'\), an approximation to it, nor any relation between \(a'\) and
\(a^{-1}\).

Combining (14), (15), and (21) proves the new arithmetic lower bound \[
\frac{\log|\Delta_H|}{M_HH}
\ge-2(1-b_H)-2E_{\rm ar}-U-o(1).                  \tag{22}
\]

\subsection{6.1 Conservative coefficient bound used by the Lean
development}\label{conservative-coefficient-bound-used-by-the-lean-development}

The module \path{ConjugateBound.lean} instead uses the convenient
one-variable estimate \(|G_i(t)|\le2\), obtained from the inherited
bounds for \(\log(1+t)\) and its tail. Bounding
\(\binom\gamma\beta\le2^{|\gamma|}\) then replaces \(KR+2\) in (20) by
\(2(KR+2)\). This proves a slightly weaker bound
\(\widetilde U=U+\log2/w_*\). For exact agreement with that
implementation, replace \(\log(KR+2)\) in (10) and the total error below
by \(\log(2(KR+2))\), and use \(\widetilde U\) in (21)--(22). All
parameter choices and the final argument are unchanged: the additional
fixed numerator is chosen before \(w_*\). This variant avoids a
multivariate Cauchy formalization and still gives exponent four.

\subsection{6.2 A fixed bound for every
embedding}\label{a-fixed-bound-for-every-embedding}

The formal parameter construction uses the slightly larger constant \[
\kappa=\sum_{\sigma:L\hookrightarrow\mathbb C}\log\max(1,|\sigma(a)|)
\] in place of \(L'\). Every summand is nonnegative, so this constant
bounds the required other-embedding cost. It is fixed before the
dimension is chosen, and hence the same dimension limit absorbs it. The
norm argument runs over all complex embeddings of the actual field
\(\mathbb Q(a)\). Since its degree is two, removing the distinguished
embedding leaves exactly one term. This avoids choosing a named
conjugate or assuming its sign, while preserving the exact coefficient
two on the rational clearing factor.

\section{7. Contradiction and
endpoint}\label{contradiction-and-endpoint}

Let \(E=2E_{\rm ar}+U+E_{\rm an}\). Expanding (15)--(17) and (20) yields
exactly \[
\begin{aligned}
E={}&\frac\nu F+
\frac{2\Lambda Fm+3\log2}{v_0}
+\frac{(\rho+2\delta+L')K+\log(3/2)}{w_0}
\\
&+2\Lambda\sum_i\frac1{w_i}
+\frac{2\theta+\log4+\log(2K)+\nu+\log(2\rho K)+\log(KR+2)}{w_*}.
\end{aligned}
\] By our choices, \(E<\varepsilon<\min(g,1)\) and \(c_\infty>3>2+E\).
Comparing (18) and (22) forces at least one of \[
c_H\le2(1-b_H)+E+o(1),                             \tag{23}
\] \[
\nu(A(1-\eta)-b_H)\le2(1-b_H)+E+o(1).              \tag{24}
\] The first is impossible for all sufficiently large \(H\), since
\(b_H\ge0\) and \(c_H\to c_\infty>2+E\). For the second, \[
\nu(A(1-\eta)-b_H)-2(1-b_H)
=g+(\nu-2)(\theta-b_H)\ge g>E,
\] so it too is impossible for all sufficiently large \(H\). The
\(o(1)\) terms are uniform in the selected columns and \(b_H\); take a
fixed margin smaller than both strict gaps. Interpolation supplies
arbitrarily large admissible \(H\), completing the contradiction.

Thus (1) holds. Moreover \(\xi\) is irrational: if \(\xi=u/v\), the
unreduced fractions \(ku/(kv)\) violate (1) for arbitrarily large
denominators. Therefore the usual exponent is defined on an irrational
real and is at most four. Its defining exponent set is nonempty by
Dirichlet and bounded above by four; if using a real supremum in Lean,
those facts must be supplied explicitly. We have not proved the endpoint
estimate with exponent exactly four and a fixed constant.

\section{8. Nonvacuity and checked
declarations}\label{nonvacuity-and-checked-declarations}

The first example is \[
a_1=(1+\sqrt2)/2,\quad a_1'=(1-\sqrt2)/2,\quad N(a_1)=-1/4.
\] It is positive, not one, quadratic, and nonintegral; \(d=2\) gives
\(2a_1=1+\sqrt2\), an algebraic integer. Here \(|a_1'|<1\), so \(L'=0\).
A stronger arithmetic-cost test is \[
a_2=(1+3\sqrt2)/2,\quad a_2'=(1-3\sqrt2)/2,\quad N(a_2)=-17/4.
\] Again \(d=2\) works, while \(a_2'<-1\) and
\(L'=\log((3\sqrt2-1)/2)>0\). Neither test has norm one; the second
exercises the nonzero conjugate-growth cost. Their integer polynomials
are \(4X^2-4X-1\) and \(4X^2-4X-17\), respectively. Irreducibility
follows from irrationality of \(\sqrt2\).

The accompanying source directory is \path{lean/AlgebraicLog/}. All
names in the following list have prefix \path{OAI.AlgebraicLog}.

\begin{itemize}
\tightlist
\item
  Ordinary integer coefficient clearing:
  \path{IntegerClearing.entry_truncatedLog_cleared_int}
\item
  One integral minor and the same determinant at every embedding:
  \path{MatrixArithmetic.exists_integral_minor_det_lift}
\item
  Actual quadratic norm inequality:
  \path{QuadraticNorm.quadratic_literal_minor_log_bound}
\item
  Normalized arithmetic bound with the row saving doubled:
  \path{SelectedArithmetic.normalized_selected_lower}
\item
  Actual conjugate determinant bound and vanishing remainder:
  \path{ConjugateAggregate.complexMinor_log_bound},
  \path{tendsto_conjugateRemainder}
\item
  Full numerical parameters from actual-\(\nu\) approximants:
  \path{exists_admissible_parameters}
\item
  Distinguished analytic bound:
  \path{LiteralAnalytic.actual_minor_analytic_bound}
\item
  Cofinal interpolation of the actual matrix:
  \path{GlobalMatrixInterpolation.cofinal_actualMatrix}
\item
  Contradiction for actual algebraic data:
  \path{DeterminantContradiction.no_fixed_data}
\item
  Actual rational adjunction and integer clearing:
  \path{FieldRealization.realField_finrank},
  \path{exists_positive_integral_multiple}
\item
  Literal quantified endpoint (1): \path{quadratic_log_explicit_bound}
\item
  Conventional irrationality-exponent consequence:
  \path{quadratic_log_irrationalityExponent_le_four}
\item
  Both nontrivial test inputs: \path{Examples.first_sample_bound},
  \path{Examples.second_sample_bound}
\end{itemize}

The final endpoint uses the standard hypotheses
\texttt{IsAlgebraic Q alpha} and
\texttt{(minpoly Q alpha).natDegree = 2}, not a custom data record. The
parameter record is constructed inside the proof from a hypothetical
failure of (1). The real-supremum consequence proves irrationality from
the unreduced lower bound, uses Dirichlet to supply a nonempty exponent
set, and proves that every member is at most four. The two examples'
algebraicity and degree are also proved, rather than assumed.

The accompanying audit prints the final statements and axiom lists and
checks the explicit target proposition. The verification record
describes native builds, exported-proof checks, and automated reviews,
together with their dependency trust boundaries. See
\href{VERIFY.md}{verification and reproduction}.

No invariant has been transported through multiplication by an
irrational algebraic number. No norm-one theorem, Galois averaging of
approximation estimates, or auxiliary field extension is used.
Mathematical novelty and external human verification remain separate
questions from this development.

\section{AI use}\label{ai-use}

The author initiated and directed an investigation of extending OpenAI's
irrationality-exponent method to logarithms of positive quadratic
algebraic numbers. OpenAI coding agents contributed substantially to
developing the candidate argument, producing the Lean formalization, and
conducting automated adversarial reviews. The construction builds on
OpenAI's published manuscript and released Lean library, the preceding
rational-logarithm extension, and Mathlib. Automated checking and model
review are distinct from independent human mathematical review.

\section{References and attribution}\label{references-and-attribution}

\begin{enumerate}
\def\labelenumi{\arabic{enumi}.}
\tightlist
\item
  OpenAI, \emph{The irrationality exponent of \(\pi\) is 2}, manuscript
  dated September 24, 2026, and
  \href{https://github.com/openai/math/tree/adc7f1241b42e322a6451854ab7e4b4c146bf78a/preprints/The-irrationality-exponent-of-pi-is-2-September-24-2026}{released
  \(\pi\) library}. The weighted interpolation and analytic determinant
  framework comes from this development.
\item
  Ryan Matthew Casper, \emph{The irrationality exponent of positive
  rational logarithms is 2}, October 7, 2026,
  \href{https://github.com/Mattie/math/tree/0cfe10002ea95c45721b93d18bffbe2c3bbbcbbd/preprints/The-irrationality-exponent-of-positive-rational-logarithms-is-2-October-7-2026}{pinned
  manuscript and formalization}. This supplies the inherited
  moving-centre interpolation and logarithmic matrix interfaces used
  here.
\item
  The mathlib Community,
  \href{https://github.com/leanprover-community/mathlib4/tree/d13f23b723b8a846827a245b89c10fc7d3f11612}{Mathlib},
  pinned revision used by the formalization.
\end{enumerate}

The proposed contribution is the general quadratic upper bound and its
formal extension. Irrationality of the logarithm of a nontrivial
positive algebraic number is classical. This paper makes no exhaustive
novelty or priority claim.
\end{document}
